\section{The pointwise count}
Let \(\A=\{a\in\Z_{>0}:p\mid a\Rightarrow p\equiv1\bmod3\}\), where \(p\)
always denotes a prime. Write \(3\!\cdot\!\A=\{3a:a\in\A\}\).
We count ordered pairs \((s_1,s_2)\) of positive integers summing to
\(n\), with
\[
(s_1,s_2)\in
\begin{cases}
\A\times\A,&n\equiv2\pmod6,\\
\A\times(3\!\cdot\!\A)\text{ or }(3\!\cdot\!\A)\times\A,
 &n\equiv4\pmod6.
\end{cases}
\]
In the second line one orientation is prescribed in advance.
For \(0<x<n\) let \(r(n;x,H)\) count these pairs with
\(|s_1-x|\le H/2\). In particular, if the first summand is \(3a\),
it is \(3a\) that lies in the interval.

\begin{theorem}\label{thm:main}
Fix \(0<\varepsilon<1\) and put \(M=\min(x,n-x)\). Uniformly for
every sufficiently large admissible \(n\) and every real \(x,H\) with
\[
0<x<n,\qquad 1\le H\le M,\qquad H^3\ge n^{1+\varepsilon}M,
\]
one has, in each of the prescribed families,
\begin{equation}\label{eq:main}
r(n;x,H)\asymp_\varepsilon \frac{HJ(n)}{\log n},
\qquad
J(n)=\prod_{\substack{p\mid n,\ p\ge5\\p\equiv2\ (3)}}
             \frac{p-1}{p-2}.
\end{equation}
The constants in the two bounds and the threshold for \(n\) depend
only on \(\varepsilon\).
\end{theorem}

The product has a direct interpretation. At \(p=5\), for example,
the two conditions \(5\nmid s_1\) and \(5\nmid n-s_1\) forbid two
residues when \(5\nmid n\), and one when \(5\mid n\). The ratio of
the allowed proportions is \((4/5)/(3/5)=4/3\). The same calculation
explains every factor of \(J(n)\) and yields an upper bound by an
ordinary sieve on integers.

The lower bound begins with a larger family that can be counted
spectrally. Write \(N(z)=z\bar z\) for the Eisenstein norm on
\(\Z[\omega]\), with \(\omega^2+\omega+1=0\). The equation
\(N(z)+N(w)=n\) is a quaternion norm equation. At an inert prime
\(p\equiv2\bmod3\), \(p\mid N(z)\) forces \(z\in p\Z[\omega]\).
It is therefore possible to remove the large inert primes by a
lattice count within the prescribed window.

A separate issue is that many coordinate pairs can give the same
two norm values. Sifting out small split primes bounds this
multiplicity, with a baseline coordinate count of size
\(H/(\log n)^2\) before the target factors. To gain a factor of \(\log n\)
we allow small split primes to occur to exponent one in one
summand, or in both summands when they divide the target.
For instance, \(7\nmid n\) permits placing \(7\) in either summand.
If \(7\mid n\), its use as a common factor reduces the target to
\(n/7\). The split-prime local factors for these two situations
supply the same logarithmic gain. The prescribed factors can
be read off the resulting values, and each pair then has bounded
weighted coordinate multiplicity.

The angular estimate costs \(\sqrt{M/H}\) in the Hecke remainder.
Multiplying by the square-root error from the norm shell gives
\(\sqrt{nM/H}\). Comparing this with the main-term width \(H\)
explains \(H^3\ge n^{1+\varepsilon}M\). The proof below treats
the finite level and the angular variable together, which lets
us prescribe the mixed orientation at any center.

\paragraph{Relation to earlier work.}
Dietmann and Elsholtz \cite{DietmannElsholtz} prove eventual
binary existence for the Gaussian prime restriction \(p\equiv1\bmod4\)
and obtain \(\gg_\delta n^{1-\delta}\) pairs for every \(\delta>0\)
without restricting the summand sizes. Howe, Huffman, Kohlmeyer and
Kruer \cite[Theorem~1.1]{Euler} prove existence for every positive
target \(2\bmod4\), using quaternionic norm counts and a sieve on the
representing vectors. The prime sets used here and in the Gaussian
problem are incomparable, as \(7\) and \(17\) show.
We use the same norm-counting principle for the Eisenstein ring.
The new estimates needed here are angular localization with prescribed
residue conditions, the densities when a sieving prime divides the
target, and a lower bound that counts distinct value pairs at the
integer-sieve scale. Ramar\'e and Viswanadham
\cite[Introduction and Theorem~1.6]{RV} discuss the split-prime
Eisenstein set in their sieve framework and prove a ternary
additive theorem for primitive Gaussian values.

\begin{corollary}\label{cor:ranges}
The count in every fixed interval of positive length for \(s_1/n\)
inside \((0,1)\) is of order \(nJ(n)/\log n\). In particular the
unrestricted count has this order. At any fixed interior share,
windows of width \(n^\theta\) have count of order
\(n^\theta J(n)/\log n\) for \(2/3<\theta<1\).
For \(1/2<\theta<1\) and fixed \(0\le a<b\), the same count holds
with \(s_1\in[a n^\theta,b n^\theta]\).
\end{corollary}

The results give two-sided order bounds. A leading constant and
a numerical onset are open here, as are the central and
small-summand endpoint exponents \(2/3\) and \(1/2\).
